\section{ความเป็นมาและความสำคัญของปัญหา}

การลดลงของกิจกรรมทางกายในชีวิตประจำวันจากการพึ่งพาเทคโนโลยีและการใช้ยานพาหนะเป็นหลัก เป็นปัญหาสุขภาพระดับโลก การเปลี่ยนแปลงพฤติกรรมเหล่านี้สัมพันธ์โดยตรงกับการเพิ่มขึ้นของอัตราโรคอ้วนและโรคไม่ติดต่อเรื้อรัง (NCDs) เช่น โรคเบาหวาน โรคหัวใจ และมะเร็ง ซึ่งส่งผลให้อัตราการป่วยและเสียชีวิตสูงขึ้น พร้อมกับสร้างภาระทางเศรษฐกิจต่อระบบสาธารณสุขของประเทศ องค์การอนามัยโลกชี้ว่าการไม่มีกิจกรรมทางกายเป็นสาเหตุการตายอันดับที่ 4 ของโลกและเป็นปัจจัยเสี่ยงหลักของโรคไม่ติดต่อเรื้อรัง ส่งผลกระทบต่อคุณภาพชีวิตและเพิ่มค่าใช้จ่ายทางการแพทย์อย่างมีนัยสำคัญ \cite{dhuli2022}

การวิ่งเป็นกิจกรรมทางกายที่เข้าถึงง่าย ไม่ต้องใช้อุปกรณ์ซับซ้อน และมีประสิทธิภาพในการป้องกันโรคไม่ติดต่อเรื้อรัง อย่างไรก็ตาม ผู้ที่เริ่มวิ่งจำนวนมากมักหยุดออกกำลังกายในระยะเวลาอันสั้น เนื่องจากขาดแรงจูงใจที่เหมาะสมและระบบสนับสนุนอย่างต่อเนื่อง ปัญหาหลักประกอบด้วยการขาดเป้าหมายที่ชัดเจน การไม่มีระบบติดตามความก้าวหน้า และการขาดปฏิสัมพันธ์เชิงบวกจากสังคมรอบข้าง งานวิจัยพบว่าแรงจูงใจมีบทบาทสำคัญต่อการรักษาพฤติกรรมการออกกำลังกายให้ยั่งยืน โดยการสนับสนุนจากผู้อื่น การออกแบบกิจกรรมที่สนุกสนาน การให้โอกาสเลือก และการสร้างความรู้สึกมีคุณค่าในตนเอง สามารถเพิ่มความตั้งใจและความถี่ในการออกกำลังกายได้อย่างมีนัยสำคัญ \cite{cachon2023}

ด้วยเหตุนี้ แนวคิดเกมมิฟิเคชัน (Gamification) จึงถูกนำมาประยุกต์ใช้ในการส่งเสริมกิจกรรมทางกาย โดยใช้องค์ประกอบของเกม เช่น คะแนน เหรียญตรา กระดานผู้นำ และภารกิจ เพื่อกระตุ้นแรงจูงใจและการมีส่วนร่วมของผู้ใช้งาน \cite{fu2024} อย่างไรก็ตาม การประยุกต์ใช้เกมมิฟิเคชันเพียงอย่างเดียวอาจมีข้อจำกัด หากขาดการบูรณาการระบบสังคม (Social System) ซึ่งเป็นองค์ประกอบสำคัญในการสร้างแรงจูงใจผ่านการปฏิสัมพันธ์และการสนับสนุนซึ่งกันและกัน งานวิจัยยืนยันว่าการเป็นส่วนหนึ่งของกลุ่มออกกำลังกายมีความสัมพันธ์กับการได้รับการสนับสนุนทางสังคม การสร้างอัตลักษณ์ของการเป็นนักออกกำลังกาย และการเพิ่มปริมาณกิจกรรมทางกาย \cite{golaszewski2022} อีกทั้งแอปพลิเคชันสุขภาพบนมือถือที่ผสมผสานองค์ประกอบทางสังคม เช่น การสร้างทีม การแข่งขันระหว่างกลุ่ม และการแบ่งปันความสำเร็จ มีประสิทธิภาพในการเพิ่มระดับกิจกรรมทางกายและรักษาพฤติกรรมในระยะยาว \cite{tong2018}

แนวทางดังกล่าวสอดคล้องกับทฤษฎีการกำหนดใจตนเอง (Self-Determination Theory: SDT) ซึ่งระบุว่าการส่งเสริมให้บุคคลรู้สึกถึงความเป็นอิสระ (Autonomy) ความสามารถ (Competence) และความสัมพันธ์ทางสังคม (Relatedness) เป็นหัวใจสำคัญของการสร้างแรงจูงใจที่ยั่งยืน แอปพลิเคชันที่ออกแบบโดยคำนึงถึงองค์ประกอบทั้งสามนี้ช่วยเพิ่มการมีส่วนร่วมและส่งผลให้เกิดการเปลี่ยนแปลงพฤติกรรมสุขภาพที่คงทนกว่า \cite{ryan2020}

ดังนั้น คณะผู้จัดทำจึงพัฒนาแอปพลิเคชัน \textit{Pacegasus} ซึ่งผสานหลักการเกมมิฟิเคชันเข้ากับระบบกลุ่มเพื่อนและเอนจินปรับแผนฝึกอัตโนมัติ เพื่อสร้างสภาพแวดล้อมที่เอื้อต่อการแข่งขันเชิงบวก การสนับสนุนทางสังคม และการแบ่งปันประสบการณ์ อันจะช่วยสร้างแรงจูงใจทั้งจากภายในและภายนอก ลดภาระโรคไม่ติดต่อเรื้อรัง ยกระดับคุณภาพชีวิตของประชาชน และลดค่าใช้จ่ายด้านสาธารณสุขของประเทศในระยะยาว

\section{วัตถุประสงค์}

\begin{enumerate}[label=\arabic*)]
    \item เพื่อออกแบบและพัฒนาแอปพลิเคชันบนมือถือสำหรับการเดินและวิ่งออกกำลังกาย ที่สามารถติดตามและบันทึกข้อมูลการเคลื่อนไหวของผู้ใช้
    \item เพื่อพัฒนาเอนจินปรับแผนฝึกอัตโนมัติ โดยใช้ข้อมูลระยะทางจริง เพซ และความสม่ำเสมอ เพื่อปรับภารกิจรายวันและรายสัปดาห์สำหรับผู้ใช้แต่ละคน
    \item เพื่อพัฒนาแอปพลิเคชันบนมือถือที่เพิ่มความน่าสนใจและแรงจูงใจในการออกกำลังกาย โดยผสานแนวคิดเกมมิฟิเคชันและการมีปฏิสัมพันธ์ทางสังคม
\end{enumerate}

\section{ขอบเขตและข้อจำกัดของระบบ}

\subsection{ขอบเขต}

โครงงานนี้ครอบคลุมการออกแบบและพัฒนาแอปพลิเคชันบนมือถือ \textit{Pacegasus} ที่มีระบบควบคุมสิทธิ์การเข้าถึง โดยกำหนดขอบเขตการใช้งานตามบทบาทของผู้ใช้งาน ดังนี้

\begin{enumerate}[label=\arabic*.\arabic*)]
    \item \textbf{ผู้ใช้ทั่วไปที่ยังไม่ได้สมัครสมาชิก (Guest)} สามารถเข้าถึงหน้าแอปพลิเคชันสาธารณะและหน้าเข้าสู่ระบบได้ แต่ไม่มีสิทธิ์ใช้งานฟีเจอร์หลัก เช่น การติดตามการวิ่ง การรับภารกิจ ระบบเกมมิฟิเคชัน หรือชุมชนนักวิ่ง และไม่สามารถจัดการหรือเข้าถึงข้อมูลภายในระบบได้
    \item \textbf{ผู้ใช้งานหรือนักวิ่ง (User/Runner)} สามารถสมัครสมาชิกและเข้าสู่ระบบด้วยชื่อผู้ใช้หรืออีเมล ติดตามและบันทึกข้อมูลการวิ่งหรือเดิน ได้แก่ ระยะทาง เพซ เวลา และแคลอรีที่เผาผลาญผ่าน GPS และเซนเซอร์ของโทรศัพท์มือถือ รับแผนการฝึกส่วนบุคคลที่ปรับอัตโนมัติตามพฤติกรรมและความก้าวหน้า เข้าร่วมระบบเกมมิฟิเคชัน เช่น การสะสมคะแนน เหรียญตรา การเลื่อนระดับ ภารกิจ และ Story Mode ใช้งานฟีเจอร์ทางสังคม เช่น การเพิ่มเพื่อน การส่งกำลังใจ การสร้างหรือเข้าร่วมกิจกรรมวิ่งกลุ่ม การดูกระดานจัดอันดับ การแชร์ผลการวิ่งไปยังสื่อสังคมออนไลน์ และจัดการข้อมูลโปรไฟล์ของตนเองได้
    \item \textbf{ผู้ดูแลระบบ (Admin)} สามารถเข้าสู่ระบบในฐานะผู้ดูแล ดูและแก้ไขข้อมูลของผู้ใช้งานทุกประเภท รวมถึงโปรไฟล์และประวัติการวิ่ง จัดการกิจกรรมและภารกิจ ตรวจสอบสถิติการใช้งานและพฤติกรรมของผู้ใช้ในภาพรวม ตลอดจนจัดการบัญชีผู้ใช้และสิทธิ์การเข้าถึงระบบ
\end{enumerate}

\subsection{ข้อจำกัดของโครงงาน}

แอปพลิเคชันต้นแบบมุ่งพัฒนาสำหรับอุปกรณ์ Android และใช้เซนเซอร์ GPS กับ Accelerometer ที่มีอยู่ในสมาร์ตโฟน การประเมินระบบเป็นการประเมินต้นแบบภายในขอบเขตของโครงงาน จึงยังไม่ครอบคลุมการใช้งานในประชากรขนาดใหญ่หรือการติดตามผลระยะยาว

\section{ประโยชน์ที่คาดว่าจะได้รับ}

\begin{enumerate}[label=\arabic*)]
    \item ได้ต้นแบบแอปพลิเคชัน \textit{Pacegasus} ที่ใช้งานได้จริง สามารถติดตามและบันทึกข้อมูลการออกกำลังกาย เช่น ระยะทาง ความเร็ว และเส้นทาง ด้วย GPS และเซนเซอร์บนโทรศัพท์มือถือ เพื่อประเมินผลและพัฒนาสมรรถภาพของผู้ใช้
    \item ได้ระบบปรับแผนฝึกอัตโนมัติ (Adaptive Training Engine) ที่สามารถสร้างและปรับเปลี่ยนภารกิจการออกกำลังกายรายวันหรือรายสัปดาห์ให้เหมาะสมกับศักยภาพและความก้าวหน้าของผู้ใช้แต่ละบุคคล เพื่อส่งเสริมการฝึกที่ปลอดภัยและมีประสิทธิภาพ
    \item เพิ่มแรงจูงใจและความสม่ำเสมอในการออกกำลังกายผ่านองค์ประกอบเกมมิฟิเคชันและการสร้างปฏิสัมพันธ์ทางสังคม เช่น การสะสมคะแนน การเข้าร่วมกิจกรรมพิเศษ และการเชื่อมต่อกับเพื่อนในแอปพลิเคชัน
    \item เป็นแนวทางในการพัฒนานวัตกรรมด้านสุขภาพดิจิทัล โดยนำเสนอต้นแบบสถาปัตยกรรมและกรอบแนวคิดโค้ชดิจิทัลเชิงปรับตัว (Adaptive Digital Coach) ที่สามารถนำไปต่อยอดเชิงวิชาการ การเรียนการสอน หรือเป็นเครื่องมือส่งเสริมกิจกรรมทางกายสำหรับบุคคลและองค์กรได้
\end{enumerate}
